% Generated by scripts/generate_future_work_status.py; do not edit.
v2実機反復・操作PSP上界 & 限定完了 & 同一UF2・同一決定的入力を2 bootで完走し、transcriptとPSP深さが一致。測定ELFとtext一致するclean再構築版で操作PSP上界を92192/120664/20980 bytesとし、全て128 KiB予約内。handler/MSPを含む全program上界ではない。 \\
v3 D1複数入力・複数配置 & 限定完了 & 10公式ベクトルを2実装×2配置で検査。正当K/S/V 40件と改変拒否40件が通り、1024-byte配置移動後もPSP深さの不一致は0/80。 \\
v3 D3二関数・全parameter & 限定完了 & 二関数の三領域を重畳。公式版/D3の全三parameterでAPI・self-test・各100 responseを通過。異なるparameter/request組は300組、二実装で計600回の検査となり、Arm frameは全6比較で減少。RP2350 p324\_3の一経路ではSign PSPを4664 bytes削減。他二parameterのwhole-image実機fitとD3時間分布は未評価。 \\
v3 stack上界 & 限定完了 & 固定frame試作で動的frame記録を19件から0件へ変更したが、実測PSPは減らなかった。link済みK/S/V根のsoftware callと最大例外entryを含むPSP上界は108300/127932/40468 bytesで、全て128 KiB予約内。clean closure imageの公式vector 0でも全観測PSPが上界内。割込み候補4根の直接call metadataは閉じたが、handler側間接callback 18地点とIRQ/MSP nestingが残るため全program上界ではない。 \\
配置・図表の自動生成 & 試作完了 & contract compilerが172080 bytesを再現し、宣言memberと直接accessのcoverage gateを追加。証拠生成器がSRAM・flash・cycle関係とPareto座標を再計算する。phase、間接alias、escape、消去の完全性は人手監査に依存する。 \\
v3公開後の文献再監査 & 完了 & 2026-09-04時点の検索式・データベース・除外基準・最近接反例を保存。否定的検索結果はv2の六条件の論理積だけに限定した。 \\
v2固定work・漏洩再検査 & 実験完了（負） & 固定budget samplerとCornacchia固定scheduleを統合した。全Signの制御・address差は再現し、残存面12件が開いているため耐性は不成立。 \\
v3漏洩評価 & 実験完了（負） & targetの10入力順位に加え、hostでmessage・署名乱数を固定した10鍵×30反復×2順序を実施し、公式版とv3適応版の双方で鍵別順位を再現した。RP2350でも同じ10鍵・固定message/RNGの200 Signを測定し、二順序で鍵対応timingを検出した。host構造traceを各実装2 processで実行し、同一鍵control 8組は一致。鍵対応の制御差=true、address差=true。物理漏洩は未判定。 \\
電力・EM実験 & 外部機材待ち & 取得firmware・解析器・合成positive controlは準備済み。電流/EM probeとscope/SCA装置が未接続で、物理traceは0件。 \\
他PQC・複数マイコン & 対象外 & 著者方針に従い本改訂では実施しない。一般性の主張にも用いない。
